
\section{例子}

下面的这些例子按专题组织，而不按难度组织。

\textbf{21.3.1. 曲线的\trGeometricGenus。}\quad 一条（域 $k$ 上的）\trSmooth{}既约射影\trCurve{} $C$ 具有\textbf{\trGeometricGenus}（geometric genus）$\mathrm{h}^0(C,\Omega_{C/k})$。（这将推广到更高的维数，见 \S21.5.3。）它总是有限的，因为 $\Omega_{C/k}$ 是\trCoherentSheaf（\mbox{练习~21.2.P}），而射影 $k$-\trScheme{}上\trCoherentSheaf{}的\trSection{}空间是有限维的（\mbox{定理~18.1.3(i)}）。可叹的是，由下述练习可见，这其实算不上一个新不变量。

\textbf{21.3.A.（容易的练习）}\quad 假设\mbox{奇迹~18.5.2}（即 $\Omega_{C/k}$ 是\trSerre{}对偶化的），证明：对一条\trSmooth{}既约射影\trCurve{} $C$ 而言，\trGeometricGenus{}等于算术亏格，即 $\mathrm{h}^0(C,\Omega_{C/k})=\mathrm{h}^1(C,\OO_C)$。

现在我们来计算一些微分吧！

\textbf{21.3.2. 射影直线。}\quad 作为一个重要的首个例子，考虑 $\mathbb{P}^1_k$，带有通常的射影\trCoordinate{} $x_0$ 与 $x_1$。一如往常，第一张片对应于 $x_0\ne0$，形如 $\Spec k[x_{1/0}]$，其中 $x_{1/0}=x_1/x_0$。第二张片对应于 $x_1\ne0$，形如 $\Spec k[x_{0/1}]$，其中 $x_{0/1}=x_0/x_1$。

两张片都同构于 $\mathbb{A}^1_k$，且 $\Omega_{\mathbb{A}^1_k}\cong\OO_{\mathbb{A}^1_k}$。（更精确地说，$\Omega_{k[x]/k}=k[x]\,\mathrm{d}x$。）于是 $\Omega_{\mathbb{P}^1_k}$ 是一个\trInvertibleSheaf（即\trLineBundle）。$\mathbb{P}^1_k$ 上的诸\trInvertibleSheaf{}皆形如 $\OO(m)$。那么，$\Omega_{\mathbb{P}^1/k}$ 究竟是哪一个\trInvertibleSheaf{}呢？

让我们取第一张片上的一个\trSection，即 $\mathrm{d}x_{1/0}$。它在上面既无\trZeros{}也无\trPole，所以让我们看看它在另一张片上的情形。由于 $x_{1/0}=1/x_{0/1}$，我们有 $\mathrm{d}x_{1/0}=-(1/x_{0/1}^2)\,\mathrm{d}x_{0/1}$。于是这个\trSection{}在 $x_{0/1}=0$ 处有一个二阶极点。故 $\Omega_{\mathbb{P}^1/k}\cong\OO(-2)$。作为推论，我们便证明了 $\mathbb{P}^1_k$ 的\trGeometricGenus{}为 $0$。

注意，若把 $k$ 换成 $\mathbb{Z}$，上面的论证照样奏效：第 15 章的 \trWeilDivisor{}与线丛理论在此适用（$\mathbb{P}^1_\mathbb{Z}$ 是\trFactorial{}的），于是前面的论证几乎不加改动地表明 $\Omega_{\mathbb{P}^1_\mathbb{Z}/\mathbb{Z}}\cong\OO(-2)$。又由于 $\Omega$ 对\trBaseChange{}行为良好（\mbox{定理~21.2.26(b)}），且任何\trScheme{}都映射到 $\Spec\mathbb{Z}$，这便蕴涵：对\emph{任何}基概形 $B$，都有 $\Omega_{\mathbb{P}^1_B/B}\cong\OO_{\mathbb{P}^1_B}(-2)$。

（另外，正如 \S18.5.2 所提示的，这表明 $\Omega_{\mathbb{P}^1/k}$ 是 $\mathbb{P}^1_k$ 的\trCanonical{}层；另见\mbox{例~18.5.4}。不过，鉴于我们尚未证明\trSerreDuality，这一点目前意义不大。）

题外一提：切丛的次数为 $2$ 这一事实，与「毛球定理」（Hairy Ball Theorem）有关（即 [Hat, Thm.~2.28] 的维数 $2$ 情形）。

\textbf{21.3.3. 超椭圆曲线。}\quad 在这一整段关于超椭圆曲线的讨论中，我们假定 $k=\bar{k}$ 且 $\operatorname{char}k\ne2$，从而可以应用 \S19.5 的讨论。考虑由一条正则射影\trCurve{} $C$ 给出的\trDoubleCover{} $\pi\colon C\to\mathbb{P}^1_k$，它在 $2g+2$ 个互不相同的点处\trRamification。我们将使用 \eqref{eq:hyperelliptic-two-affine} 中超椭圆曲线的显式坐标描述。特别地，$\pi$ 在 $0$ 处不\trRamification。由\mbox{定理~19.5.1}，$C$ 的\trGenus{}为 $g$。

\textbf{21.3.B.（练习：超椭圆曲线上的微分）}\quad 可逆层 $\Omega_{C/k}$ 的次数是多少？（提示：取 $x$ 为 $\mathbb{P}^1_k$ 的某张\trCoordinatePatch{}上的一个\trCoordinate。考虑 $C$ 上的 $\pi^*\,\mathrm{d}x$，数一数它的极点与零点。用 \S19.5 的显式坐标。你应当发现：$\pi^*\,\mathrm{d}x$ 有 $2g+2$ 个零点与四个极点，以重数计数，总计为 $2g-2$。）做完这道练习，你就能很好地为 \S21.4 的\trRiemann--\trHurwitz{}公式作好准备。

\textbf{21.3.C.（练习）}\quad 证明 $\mathrm{h}^0(C,\Omega_{C/k})=g$（从而 $C$ 的\trGeometricGenus{}为 $g$），做法如下。

（a）证明：$\Spec k[x,y]/(y^2-f(x))$ 上任一在\trInvolution{} $y\mapsto -y$ 下保持的\trRegular{}\trDifferential{} $\omega$（即 $\Omega_{(k[x,y]/(y^2-f(x)))/k}$ 的一个元素）都是从 $\Spec k[x]$ 上的某个\trDifferential{}拉回的。（提示：为陈述「$\omega/\mathrm{d}x$ 是 $\Spec k[x,y]/(y^2-f(x))$ 上一个在\trInvolution{} $y\mapsto -y$ 下保持的\trRationalFunction」赋予意义，并证明 $\omega/\mathrm{d}x$ 是 $x$ 的\trRationalFunction。）

（b）用（a）证明：任一属于 $\mathrm{H}^0(C,\Omega_{C/k})$ 且在\trInvolution{} $\iota\colon y\mapsto -y$ 下保持的\trDifferential{} $\omega$ 都必是经由 $\pi$ 从 $\mathbb{P}^1$ 拉回的，从而必为零。证明：每个\trDifferential{} $\omega\in\mathrm{H}^0(C,\Omega_{C/k})$ 都满足 $\iota^*\omega=-\omega$。

（c）证明：$\mathrm{d}x/y$ 是 $\Spec k[x,y]/(y^2-f(x))$ 上的一个（\trRegular）\trDifferential。证明：对 $0\le i<g$，$x^i(\mathrm{d}x)/y$ 可延拓为 $C$ 上一个整体\trDifferential{} $\omega_i$（即无极点）。

（d）证明：诸 $\omega_i$（$0\le i<g$）是线性无关的\trDifferential，也就是说，它们在向量空间 $\mathrm{H}^0(C,\Omega_{C/k})$ 中线性无关。（提示：令 $\{p,q\}=\pi^{-1}(0)$。证明 $\omega_i$ 在 $p$ 与 $q$ 两处的\trValuation{}均为 $i$。若 $\omega:=\sum_{j=s}^{g-1}a_j\omega_j$ 是诸 $\omega_i$ 的一个非平凡线性组合，其中 $a_j\in k$ 而 $a_s\ne0$，证明 $\omega$ 在 $p$ 处的\trValuation{}为 $s$，从而 $\omega\ne0$。）

（e）证明：诸 $\omega_i$ 张成诸\trDifferential{}的向量空间 $\mathrm{H}^0(C,\Omega_{C/k})$。（提示：若 $\omega\in\mathrm{H}^0(C,\Omega_{C/k})$，用（d）证明存在唯一的诸 $a_i$，使 $\omega':=\omega-\sum_{i=0}^{g-1}a_i\omega_i$ 在 $p$ 处以阶 $\ge g$ 消失。由（b），$\omega'$ 在 $q$ 处也以同一阶消失。再用\mbox{练习~21.3.B} 证明 $\omega'$ 必为零。）

于是，$\Omega_{C/k}$ 是一个次数为 $2g-2$、带有 $g$ 个\trSection{}的\trInvertibleSheaf。

\textbf{21.3.D.}\ \hard\ \textbf{（练习：通向\trSerre{}对偶）}\quad （日后你或许会把它看作\trSerre{}对偶发挥作用的一个例子。）

（a）证明 $\mathrm{h}^1(C,\Omega_{C/k})=1$。把 $\mathrm{H}^1(C,\Omega_{C/k})$ 的一个生成元诠释为 $x^{-1}\,\mathrm{d}x$。（特别地，拉回映射 $\mathrm{H}^1(\mathbb{P}^1,\Omega_{\mathbb{P}^1/k})\to\mathrm{H}^1(C,\Omega_{C/k})$ 是一个\trIsomorphism。）

（b）描述一个自然的\trPerfectPairing
\[
\mathrm{H}^0(C,\Omega_{C/k})\times\mathrm{H}^1(C,\OO_C)\longrightarrow\mathrm{H}^1(C,\Omega_{C/k}).
\]

在我们显式的坐标下，你可以这样来诠释它。回忆\trHyperelliptic{}\trRiemann--\trHurwitz{}公式（\mbox{定理~19.5.1}）的证明中，$\mathrm{H}^1(C,\OO_C)$ 可被诠释为
\[
\langle y/x,\ y/x^2,\ \dots,\ y/x^g\rangle.
\]
那么，配对
\[
\langle \mathrm{d}x/y,\ x\,\mathrm{d}x/y,\ \dots,\ x^{g-1}\,\mathrm{d}x/y\rangle\times\langle y/x,\ y/x^2,\ \dots,\ y/x^g\rangle\longrightarrow\langle x^{-1}\,\mathrm{d}x\rangle
\]
基本上就是「相乘，然后读出 $x^{-1}\,\mathrm{d}x$ 那一项」。用更高档的话说：「相乘，然后取留数」。（你或许想把它与\mbox{例~18.5.4} 作个比较。）

\textbf{21.3.4. 离散赋值环。}\quad 下面这道练习将用于 \S21.4 中\trRiemann--\trHurwitz{}公式的证明。

\textbf{21.3.E.（练习）}\quad 设\trDVR{} $(A,\mathfrak{m},k)$ 是某个\trFinitelyGenerated{} $k$-代数的\trLocalization。设 $t$ 是 $A$ 的一个\trUniformizer。证明：诸\trDifferential{}是秩 $1$ 的自由模，且由 $\mathrm{d}t$ 生成，即 $\Omega_{A/k}=A\,\mathrm{d}t$，做法如下。（借 \S21.5 中的想法来证明这件事也是可能的。）

（a）证明 $\Omega_{A/k}$ 是一个\trFinitelyGenerated{} $A$-\trModule。

（b）证明 $\langle\mathrm{d}t\rangle=\Omega_{A/k}$，也就是说，$\times\,\mathrm{d}t\colon A\to\Omega_{A/k}$ 是满射，做法如下。设 $\pi$ 是投射 $A\to A/\mathfrak{m}=k$，于是对 $a\in A$，有 $a-\pi(a)\in\mathfrak{m}$。定义 $\sigma(a)=(a-\pi(a))/t$。证明 $\Omega_{A/k}=\langle\mathrm{d}t\rangle+\mathfrak{m}\Omega_{A/k}$，所用的事实是：对每个 $a\in A$，
\[
\mathrm{d}a=\sigma(a)\,\mathrm{d}t+t\,\mathrm{d}\sigma(a).
\]

把\trNakayama{}引理版本~3（\mbox{练习~8.2.G}）应用于 $\langle\mathrm{d}t\rangle\subset\Omega_{A/k}$。（这个论证几乎原封不动地也能用来证明：若 $(A,\mathfrak{m},k)$ 是 $k$ 上某个\trFinitelyGenerated{}代数的\trLocalization，而 $t_1,\dots,t_n$ 生成 $\mathfrak{m}$，则 $\mathrm{d}t_1,\dots,\mathrm{d}t_n$ 生成 $\Omega_{A/k}$。）

（c）由（b），$\Omega_{A/k}$ 是一个\trPrincipal{} $A$-\trModule。证明 $\times\,\mathrm{d}t$ 是单射，做法如下。由离散赋值环上有限生成模的分类（\mbox{注记~6.4.4}），只须证明对一切 $m$ 都有 $t^m\,\mathrm{d}t\ne0$。满射 $A\to A/(t^N)$ 诱导一个映射 $\Omega_{A/k}\to\Omega_{(A/(t^N))/k}$，所以只须证明 $t^m\,\mathrm{d}t$ 在 $\Omega_{(A/(t^N))/k}$ 中非零。证明 $A/(t^N)\cong k[t]/(t^N)$。通常的多项式求导法则给出一个映射 $k[t]/(t^N)\to k[t]/(t^{N-1})$，它是 $k[t]/(t^N)$ 在 $k$ 上的一个\trDerivation，而只要 $N$ 充分大，$t^m\,\mathrm{d}t$ 就不会被映到 $0$。把各块拼在一起，完成证明。（这些想法的延伸可以证明：若 $k$ 上某个\trFinitelyGenerated{}的代数是一个\trRegular\trLocalRing，而 $(A,\mathfrak{m},k)$ 是它的\trLocalization，则 $\Omega_{A/k}$ 是秩 $\dim A$ 的自由模。）

\textbf{21.3.5. 射影空间与欧拉正合列。}\quad 我们接下来考察射影空间 $\mathbb{P}^n_k$ 的诸\trDifferential，或者更一般地，$\mathbb{P}^n_A$ 的诸\trDifferential，其中 $A$ 是任意一个\trRing。由于射影空间由形如 $\mathbb{A}^n_A$ 的\trAffine\trOpenSubset{}所覆盖，而在这些开集上诸\trDifferential{}构成秩 $n$ 的\trLocallyFreeSheaf，所以 $\Omega_{\mathbb{P}^n_A/A}$ 也是一个秩 $n$ 的\trLocallyFreeSheaf。

\textbf{21.3.6.\ \trhTheorem（欧拉正合列）。}\ ——\ \emph{诸\trDifferential{}的层 $\Omega_{\mathbb{P}^n_A/A}$ 满足如下正合列：}
\[
0\longrightarrow\Omega_{\mathbb{P}^n_A/A}\xrightarrow{\ \alpha\ }\OO_{\mathbb{P}^n_A}(-1)^{\oplus(n+1)}\xrightarrow{\ \phi\ }\OO_{\mathbb{P}^n_A}\longrightarrow 0.
\]

这条正合列很称手，因为你可以借它以一种具体的方式把握住 $\Omega_{\mathbb{P}^n_A/A}$。应用见\mbox{练习~21.5.Q}。将这条正合列取对偶，我们便得到正合列
\[
0\longrightarrow\OO_{\mathbb{P}^n_A}\xrightarrow{\ \phi^\vee\ }\OO_{\mathbb{P}^n_A}(1)^{\oplus(n+1)}\longrightarrow\mathscr{T}_{\mathbb{P}^n_A/A}\longrightarrow 0.
\]

\textbf{21.3.7.}\ \hard\ \emph{\mbox{定理~21.3.6} 的证明。}\quad 这个证明的真正要点是：我们考虑 $\mathbb{A}^{n+1}_A\setminus\{0\}$ 上那些作为 $\mathbb{P}^n_A$ 上\trDifferential{}之拉回的\trDifferential。若想从 Koszul 复形（\S24.4.1 将简要提到）的角度另找一个解释，见 [E, \S17.5]。

我们先来描述一个映射 $\phi\colon\OO(-1)^{\oplus(n+1)}\to\OO$，稍后再把它的\trKernel{}与 $\Omega_{\mathbb{P}^n_A/A}$ 识别起来。该映射由
\[
\phi\colon(s_0,s_1,\dots,s_n)\longmapsto x_0s_0+x_1s_1+\cdots+x_ns_n.
\]
给出。（你能看出为什么吗？）你应当把它看作一个「次数 $1$」的映射，因为每个 $x_i$ 的次数都是 $1$。

\textbf{21.3.8.\ \trhRemark。}\quad $\phi$ 的\trDual{} $\phi^\vee\colon\OO\to\OO(1)^{\oplus(n+1)}$ 借\mbox{重要定理~15.2.2}（到射影空间的映射）给出一个映射 $\mathbb{P}^n\to\mathbb{P}^n$。这个映射是恒等映射，而这也是以一种「自然」（无坐标）的方式描述 $\phi$ 的途径之一。

\textbf{21.3.F.（容易的练习）}\quad 证明 $\phi$ 是满射，办法是在诸开集 $\mathrm{D}(x_i)$ 上逐一验证。（存在一行字的解法。）

现在我们必须把 $\phi$ 的\trKernel{}与诸\trDifferential{}识别起来。这件事可以在每一块 $\mathrm{D}(x_i)$ 上完成，只要我们所做的方式能对每个开集同时奏效即可。于是我们考虑开集 $\mathsf{U}_0$，即 $x_0\ne0$ 之处，我们有通常的\trCoordinate{} $x_{j/0}=x_j/x_0$（$1\le j\le n$）。给定一个\trDifferential
\[
f_1(x_{1/0},\dots,x_{n/0})\,\mathrm{d}x_{1/0}+\cdots+f_n(x_{1/0},\dots,x_{n/0})\,\mathrm{d}x_{n/0},
\]
我们必须造出 $\OO(-1)$ 的 $n+1$ 个\trSection。作为动机，我们只看第一项，并假装那些射影\trCoordinate{}就是真实的\trCoordinate。
\[
\begin{aligned}
f_1\,\mathrm{d}x_{1/0}&=f_1\,\mathrm{d}(x_1/x_0)\\
&=f_1\,\frac{x_0\,\mathrm{d}x_1-x_1\,\mathrm{d}x_0}{x_0^2}\\
&=-\frac{x_1}{x_0^2}f_1\,\mathrm{d}x_0+\frac{1}{x_0}f_1\,\mathrm{d}x_1.
\end{aligned}
\]
注意，$x_0$ 乘以「$\mathrm{d}x_0$ 的系数」再加上 $x_1$ 乘以「$\mathrm{d}x_1$ 的系数」等于 $0$，而且两个系数都是 $-1$ 次\trHomogeneous{}的。受此启发，我们定义
\begin{equation*}\tag{21.3.8.1}\label{eq:euler-alpha-map}
\alpha\colon f_1\,\mathrm{d}x_{1/0}+\cdots+f_n\,\mathrm{d}x_{n/0}\longmapsto\Bigl(-\frac{x_1}{x_0^2}f_1-\cdots-\frac{x_n}{x_0^2}f_n,\ \frac{f_1}{x_0},\ \frac{f_2}{x_0},\ \dots,\ \frac{f_n}{x_0}\Bigr).
\end{equation*}

注意，在 $\mathsf{U}_0$ 上，这的确给出了一个从 $\Omega_{\mathbb{P}^n_A/A}$ 到 $\OO(-1)^{\oplus(n+1)}$ 的单射，并且它满射到 $\OO(-1)^{\oplus(n+1)}\to\OO_X$ 的\trKernel{}之上（若 $(g_0,\dots,g_n)$ 属于该\trKernel，则对 $i>0$ 取 $f_i=x_0g_i$）。

让我们确认一下：把这个构造应用于两张不同的\trCoordinatePatch（比如说 $\mathsf{U}_0$ 与 $\mathsf{U}_1$），给出的是同一个答案。（首次阅读时最好略过这段验证。）注意
\[
\begin{aligned}
f_1\,\mathrm{d}x_{1/0}+f_2\,\mathrm{d}x_{2/0}+\cdots&=f_1\,\mathrm{d}\frac{1}{x_{0/1}}+f_2\,\mathrm{d}\frac{x_{2/1}}{x_{0/1}}+\cdots\\
&=-\frac{f_1}{x_{0/1}^2}\,\mathrm{d}x_{0/1}+\frac{f_2}{x_{0/1}}\,\mathrm{d}x_{2/1}-\frac{f_2x_{2/1}}{x_{0/1}^2}\,\mathrm{d}x_{0/1}+\cdots\\
&=-\frac{f_1+f_2x_{2/1}+\cdots}{x_{0/1}^2}\,\mathrm{d}x_{0/1}+\frac{f_2}{x_{0/1}}\,\mathrm{d}x_{2/1}+\cdots.
\end{aligned}
\]
在这个映射下，$\mathrm{d}x_{2/1}$ 这一项进到 $\OO(-1)^{\oplus(n+1)}$ 的第二个因子（这些因子的下标从 $0$ 编到 $n$），并给出 $f_2/x_0$，正合所需（对 $j>2$ 的 $\mathrm{d}x_{j/1}$ 亦类似）。另外，$\mathrm{d}x_{0/1}$ 这一项进到「零」因子，并给出
\[
-\Bigl(\sum_{i=1}^n f_i(x_i/x_1)\big/(x_0/x_1)^2\Bigr)\big/x_1=-\sum_{i=1}^n f_i(x_i/x_1)x_i\big/x_0^2,
\]
正合所需。最后，「第一个」因子必定正确，因为诸 $x_i$ 与第 $i$ 个因子的乘积对 $i$ 求和为零。

\textbf{21.3.G.（练习）}\quad 通过验证这个映射 $\alpha\colon\Omega_{\mathbb{P}^n_A/A}\to\OO_{\mathbb{P}^n_A}(-1)^{\oplus(n+1)}$ 将 $\Omega_{\mathbb{P}^n_A/A}$ 与 $\ker\phi$ 识别起来，完成\mbox{定理~21.3.6} 的证明。\hfill$\square$

\textbf{21.3.H.（\S18.5.7 中许诺的练习）}\quad 证明 $\mathrm{h}^1(\mathbb{P}^2_k,\Omega_{\mathbb{P}^2_k/k})>0$。证明 $\Omega_{\mathbb{P}^2_k/k}$ 不是诸\trLineBundle{}的\trDirectSum，甚至连把它用\trLineBundle{}\emph{滤过}都办不到。证明\mbox{定理~18.5.6}（$\mathbb{P}^1$ 上一切有限\trRank{}\trVectorBundle{}都分裂为诸\trLineBundle{}之直和）无法推广到 $\mathbb{P}^2_k$ 上秩 $2$ 的\trVectorBundle。

\textbf{21.3.9.}\ \emph{欧拉正合列的种种推广。}\quad 欧拉正合列的种种推广相当有用。我们日后并不会用到它们，所以不给证明。首先，若把 $\Spec A$ 换成任意一个基概形，这个论证可以轻而易举地推广过去。欧拉正合列还能以若干种方式进一步推广。作为第一步，设 $\mathscr{V}$ 是概形 $X$ 上秩为 $n+1$ 的一个\trLocallyFreeSheaf（即\trVectorBundle）。设 $\pi\colon\mathbb{P}\mathscr{V}\to X$ 是相应的射影丛（$\mathbb{P}\mathscr{V}$ 的定义见 \S17.2.3，其中有一个关于对偶的警示）。于是 $\Omega_{\mathbb{P}\mathscr{V}/X}$ 容于一个欧拉正合列：
\[
0\longrightarrow\Omega_{\mathbb{P}\mathscr{V}/X}\longrightarrow\OO_{\mathbb{P}\mathscr{V}}(-1)\otimes\pi^*\mathscr{V}^{\vee}\longrightarrow\OO_{\mathbb{P}\mathscr{V}}\longrightarrow 0.
\]

这件事有用并不显然，但其实我们已经在曲线的 $\mathbb{P}^1$-丛的情形里隐式地见过它：\mbox{练习~20.2.J} 中一个截面的\trNormalBundle{}正是这样被识别出来的。

\textbf{21.3.10.}\ \extrahard\ \emph{到\trGrassmannian{}的推广。}\quad 作为另一个推广，取定一个\trBaseField{} $k$，并设 $G(m,n+1)$ 是 $(n+1)$ 维\trVectorSpace{} $V$ 中由维数为 $m$ 的子向量空间构成的空间（即\trGrassmannian，\S16.4）。在 $G(m,n+1)$ 之上，我们有\trLocallyFreeSheaf{}的一个短正合列
\[
0\longrightarrow\mathscr{S}\longrightarrow\OO_{G(m,n+1)}\otimes V^{\vee}\longrightarrow\mathscr{Q}\longrightarrow 0,
\]
其中 $\OO_{G(m,n+1)}\otimes V^{\vee}$ 是「纤维皆为 $V^{\vee}$ 的\trTrivialBundle{}」（你明白这是什么意思吗？），而 $\mathscr{S}$ 是「泛向量丛」或「重言向量（子）丛」。于是，存在一个\trCanonical{}\trIsomorphism
\begin{equation*}\tag{21.3.10.1}\label{eq:grassmannian-cotangent-hom}
\Omega_{G(m,n+1)/k}\longleftrightarrow\mathscr{H}om(\mathscr{Q},\mathscr{S}).
\end{equation*}

\textbf{21.3.I.（练习）}\quad 回忆：在射影空间的情形，即 $m=1$，$\mathscr{S}=\OO(-1)$（\mbox{练习~17.1.H}）。用欧拉正合列（\mbox{定理~21.3.6}）验证此情形下的 \eqref{eq:grassmannian-cotangent-hom}。

\textbf{21.3.J.（练习）}\quad 证明 \eqref{eq:grassmannian-cotangent-hom}，并解释它如何推广\mbox{练习~20.2.I}。（\mbox{练习~20.2.I} 的提示或有所帮助。）

这个格拉斯曼流形的事实还能进一步推广到格拉斯曼丛，以及旗簇，还有旗丛。
